\chapter{Conclusions and Recommendations}
\label{ch:conclusion}

\section{Summary of Achievements}

The project delivers a complete, running data platform for the
ride-hailing fleet use case: two simulated ingestion paths, a Kafka
topic partitioned for correctness, an immutable Parquet master dataset,
a Spark Structured Streaming speed layer with stateful idle alerting, a
nine-task Airflow reconciliation over PySpark, a PostgreSQL serving store
holding both layers' views, a FastAPI surface of eleven endpoints, a
consolidated daily profitability report, two Grafana dashboards, and ten
Prometheus alert rules routed through Alertmanager. It starts from a
fresh clone with one command on a 16\,GB laptop, and 202 automated tests
plus a 28-check end-to-end smoke test pass against it.

\section{Key Findings}

The business question posed by the brief contains two sub-questions with
incompatible latency, accuracy and recomputability requirements. A Lambda
architecture is justified here not by convention but by a specific
property Kappa cannot supply economically: \textbf{restating a past day's
financial figures when a correction arrives on a different channel, days
later.}

The project demonstrates that property directly. A corrected file
triggers an idempotent recompute producing the same 50 rows with updated
values, and a forced replay reproduces bit-identical figures. It also
\emph{quantifies} the cost of the speed layer's approximation rather than
asserting it, measuring one dropped revenue-bearing event on one day of
four and explaining, arithmetically, why that is the expected magnitude.
We regard the decision to report a measured drift of \code{0.000000} on
three days, with the arithmetic that predicts it, as a stronger result
than a larger number obtained by quietly raising the injected fault rate
would have been.

Lambda's principal weakness --- divergent codebases --- is mitigated
structurally by a shared \code{common/} package, with the two
performance-motivated duplications pinned together by tests that feed
identical inputs through both implementations. This is the part of the
design we would defend most strongly to a reader who arrived persuaded by
the Kappa critique.

\section{Limitations}

The limitations are set out in full in Section~\ref{sec:limitations}. The
three that most constrain the conclusions are that batch Spark runs in
local mode, so distributed batch is designed for but not demonstrated;
that the drift figures are small integer counts from a single run, so
four days is not a sample from which to generalise a rate; and that the
simulated clock's derivation from wall-clock time makes operational
downtime indistinguishable from low demand in the resulting data.

\section{Recommendations and Future Work}

In priority order, and independent of scale:

\begin{enumerate}
  \item \textbf{Decouple the archiver from the speed layer} onto
        separate workers. They already have separate checkpoints; sharing
        cores is what turned a worker restart into 26 minutes of lost
        archiving.
  \item \textbf{Add a schema registry} in front of the telemetry topic,
        so an incompatible producer is rejected at publish time rather
        than discovered in the dead-letter queue, and give the DLQ an
        automated replay path.
  \item \textbf{Adopt a lake table format} (Iceberg or Delta Lake) to
        give the curated lake writes the same atomicity that
        delete-then-insert already gives the PostgreSQL views.
  \item \textbf{Annotate outage windows} on the simulated timeline, so
        the \code{becoming\_unprofitable} trend rule can exclude days
        degraded by pipeline downtime instead of reading them as a
        commercial decline.
  \item \textbf{Express alert thresholds as ratios} rather than absolute
        counts, so that \code{IdleVehiclesHigh} survives a change in
        fleet size.
\end{enumerate}

\section{Lessons Learned}

The most valuable outcome was not the architecture but what operating it
revealed: a green healthcheck over a dead archiver, a correct pipeline
producing economically meaningless output, and an interval that was
60$\times$ longer than intended because it was expressed in the wrong
clock. None of the three was visible in code review. Each was found by
running the system and checking the numbers against expectations, which
is the habit this project is ultimately arguing for --- and the same
habit that produced its central result, since the decision to measure why
the drift was zero, rather than to assume the measurement was broken, is
what turned an anticlimax into the report's strongest section.
